\documentclass[11pt,letterpaper]{article}
\usepackage[T1]{fontenc}
\usepackage{lmodern,microtype}
\usepackage[margin=0.85in]{geometry}
\usepackage{amsmath,amssymb,amsthm,mathtools,mathrsfs}
\usepackage{booktabs,longtable,array,enumitem}
\usepackage{xcolor,hyperref,fancyhdr}
\hypersetup{colorlinks=true,linkcolor=blue!45!black,citecolor=blue!45!black,urlcolor=blue!45!black,pdftitle={Sparse target packing and arithmetic separation for Diophantine consistency},pdfauthor={Research prepared for Vladimir Reshetnikov and ProveIt}}
\setlength{\emergencystretch}{2em}
\setlist{itemsep=3pt,topsep=5pt}
\pagestyle{fancy}\fancyhf{}
\fancyhead[L]{\small Diophantine consistency: explicit reductions}
\fancyhead[R]{\small 23 October 2026}
\fancyfoot[C]{\thepage}
\setlength{\headheight}{14pt}
\newtheorem{theorem}{Theorem}[section]
\newtheorem{lemma}[theorem]{Lemma}
\newtheorem{proposition}[theorem]{Proposition}
\newtheorem{corollary}[theorem]{Corollary}
\theoremstyle{remark}\newtheorem{remark}[theorem]{Remark}
\newcommand{\Z}{\mathbb Z}\newcommand{\N}{\mathbb N}
\newcommand{\Pos}{\mathbb Z_{>0}}
\newcommand{\Con}{\operatorname{Con}}\newcommand{\ZFC}{\mathsf{ZFC}}
\newcommand{\pc}{\operatorname{pc}}
\newcommand{\file}[1]{\nolinkurl{#1}}
\title{\vspace{-1em}\textbf{Sparse target packing and arithmetic separation\\for Diophantine consistency}\\\large Explicit equations for the absence of a ZFC refutation}
\author{Research prepared for Vladimir Reshetnikov and the ProveIt project}
\date{23 October 2026}
\begin{document}
\maketitle
\begin{abstract}
We construct explicit improvements to three measures of simplicity for a
Diophantine equation whose lack of solutions is provably equivalent, in
ZFC, to syntactic consistency of ZFC. The principal unrestricted-integer
equation has ten variables, exact total degree $54\,768\,104$, and a
division-free circuit of 709 binary arithmetic operations. A companion has
nine strictly positive variables, exact degree $33\,332\,620$, and 682
operations. Independently, a squared-modulus Pell congruence reduces the
fifteen-positive-variable operation record from 67 to 66, at exact degree
157, and gives 72-operation/degree-80 and 74-operation/degree-68 alternatives.

Three ingredients supply the improvement: coefficient extraction at
positions selected using the actual quadratic supports; translation of
positivity inside the packed word; and a square-root separator proved by
divisibility and parity instead of large magnitude bounds. Every exponent
and coefficient is fixed. The accompanying archive contains literal
circuits, coefficient payloads, deterministic builders, independent finite
certificates, source pins, and detailed dependency proofs. The new
equivalences are ordinary mathematical proofs formalizable in ZFC. We do
not claim a new end-to-end Lean/Rocq certificate, a proof of consistency,
or global minimality.
\end{abstract}
\newpage
\tableofcontents
\newpage

\section{Results, domains, and the meaning of simplicity}
The reviewed repository snapshot is
\begin{center}\small\texttt{124c07c0b77a9e20e84676514d8b1ff78536ea3f}.\end{center}
It contains the newly placed wave-6 manuscripts and the active formal
development in \file{Computability/DiophantineConsistency}. Comparisons
refer to that snapshot, rather than the superseded wave-5 summary alone.

\begin{theorem}[Explicit consistency equations]\label{thm:main}
There are the explicitly specified fixed integer polynomials $F_{10}$,
$F_9$, and $P_{66}$ in this report and archive for which ZFC proves
\begin{align*}
 \neg\exists\mathbf x\in\Z^{10}\ F_{10}(\mathbf x)=0
 &\ \longleftrightarrow\ \Con(\ZFC),\\
 \neg\exists\mathbf x\in\Pos^9\ F_9(\mathbf x)=0
 &\ \longleftrightarrow\ \Con(\ZFC),\\
 \neg\exists\mathbf x\in\Pos^{15}\ P_{66}(\mathbf x)=0
 &\ \longleftrightarrow\ \Con(\ZFC).
\end{align*}
Their exact degrees and delivered circuit sizes are in
Table~\ref{tab:new}. The theorem uses the inherited K215 calibration and
ordinary native-compiler correctness theorem. The ten-integer branch also
uses Sun's Lucas/interpolation theorems, and the nine-positive branch uses
the Jones ratio bridge. These dependencies are identified explicitly below.
\end{theorem}

\begin{table}[h]
\centering
\begin{tabular}{@{}llrr@{}}\toprule
Polynomial & Witness domain & Exact degree & Binary operations\\\midrule
$F_{10}$ & $\Z^{10}$ & $54\,768\,104$ & 709\\
$F_9$ & $\Pos^9$ & $33\,332\,620$ & 682\\
$F_{9,\N}$ & $\N^9$ & $33\,332\,620$ & 685\\
$P_{66}$ & $\Pos^{15}$ & 157 & 66\\
$P_{68}$ & $\Pos^{15}$ & 113 & 68\\
$P_{72}$ & $\Pos^{15}$ & 80 & 72\\
$P_{74}$ & $\Pos^{15}$ & 68 & 74\\\bottomrule
\end{tabular}
\caption{New explicit presentations; $\N=\{0,1,2,\ldots\}$.
Subscripts 66, 68, 72, and 74 count operations, not variables.}\label{tab:new}
\end{table}

We count every binary addition, subtraction, and multiplication,
including multiplication by a fixed numeral. Reuse is free. A fixed
integer coefficient is a leaf; its finite construction and binary value
are supplied separately. Fixed powers are expanded into multiplication
schedules. There are no division, root, sign, divisibility, or
variable-exponentiation gates.

The closest low-variable baseline is package 40: ten unrestricted
integers, degree $237\,520\,456$, and 782 operations. Package 38 has a
different ten-integer tradeoff, 773 operations at degree $449\,607\,016$.
Package 45 gives nine positive variables at degree $129\,875\,980$ and
788 operations. The new ten-integer equation improves both of these
ten-integer comparators. The nine-positive equation improves both degree
and operations against package 45. Degree drops by factors approximately
$4.337$ and $3.896$, respectively. These are concrete improvements to the
current repository; they do not settle a general minimum-variable or
minimum-degree problem.

The 66-operation equation is a different tradeoff. It raises the old P67
degree from 155 to 157. Its 72/80 and 74/68 companions improve the checked
73/80 and 75/68 operation counts without changing those degrees. The
new circuits have not themselves been submitted to a proof-assistant kernel.

\section{Semantic starting point and fixed coefficients}\label{sec:dependencies}
\subsection{The exact logical statement}
Here $\Con(\ZFC)$ means absence of a finite syntactic derivation of falsity
from the specified ZFC axiom predicate. The repository's literal machine
K215, on blank tape, is calibrated to this predicate:
\begin{equation}\label{eq:machine}
K215\text{ halts on blank input}
\quad\longleftrightarrow\quad
\text{ZFC has a finite refutation}.
\end{equation}
The formal endpoint is
\file{BusyBeaver.ZFC.KEndpoint215.blank_halts_iff_zfc_refutation}.
The explicit equations use the ordinary native-window/Pell compiler
theorem. No consistency assumption or search for a refutation enters
their coefficient construction.

For the low-variable constructions retain the \emph{complete} package-31
tuple
\[
(\beta,K_0,M_C,M_F),\qquad
\beta=B_{\rm nat}-1,\quad K_0=D_C+B_{\rm nat}D_R,
\]
where $B_{\rm nat}=2^{1064921}$ and
$M_F=M_{F,\rm native}+B_{\rm nat}-1$ is already shifted.
These are fixed integers, not quantified variables. Their binary payloads,
deterministic generator, 943-tile table, and literal K215 transition table
are bundled. Their respective bit lengths are
\[
1064921,\quad1070004,\quad1064921,\quad1064922.
\]
The short equation in Section~\ref{sec:short} uses instead the
\emph{whole} package-36 tuple with $B_{\rm nat}=2^{1012115}$.
The two tuples are never spliced or silently identified.

\subsection{The quadratic parent}
The 31 positive parent coordinates, in their data order, are
\begin{gather*}
J,F,\alpha,z,f,h,i,o,j,s,w,\tau,\eta,\zeta,y,\gamma,\\
N,E,P,c,a,X,D,C_2,L,F_2,V,g,I,B_{\rm aux},\Delta.
\end{gather*}
Write
\[
q=\beta J+1,\qquad k=\eta+\zeta,\qquad H_p=4a+3,\qquad C_p=q-F-\alpha.
\]
Table~\ref{tab:parent} gives the exact system. The last column anticipates
our new extraction positions.
\begin{longtable}{@{}r@{\quad}l r@{}}
\caption{Unchanged quadratic parent and new targets.}\label{tab:parent}\\
\toprule Row & Equation & Target\\\midrule\endfirsthead
\toprule Row & Equation & Target\\\midrule\endhead
1 & $L=q^2-1$ &714\\
2 & $N=q(L+1)$ &6707\\
3 & $X=wq$ &6262\\
4 & $a=E+sN$ &634\\
5 & $E=X(a-E)$ &6148\\
6 & $c=k(a-E)+\eta$ &5133\\
7 & $P=E(c-\eta)+k$ &4099\\
8 & $\tau^2=P(P-k)+1$ &7243\\
9 & $D=X+ac+\gamma H_p$ &1024\\
10 & $\Delta=a^2+H_p$ &2665\\
11 & $C_2=c^2$ &4617\\
12 & $D^2=\Delta C_2+1$ &5288\\
13 & $k=jc-V+1+hE$ &2271\\
14 & $jc-V=(q-F)(N-q)-C_pL+(M_C+qM_F)J$ &3114\\
15 & $(K_0+w)C_p+q-F-z(q-1)=1$ &1023\\
16 & $F_2=f^2$ &5402\\
17 & $V=of-c$ &5193\\
18 & $g=y^2-V^2$ &2781\\
19 & $I=iC_2$ &1906\\
20 & $B_{\rm aux}=I^2$ &6470\\
21 & $B_{\rm aux}=\Delta(F_2-1)$ &6168\\
22 & $y^2-B_{\rm aux}g=1$ &4539\\\bottomrule
\end{longtable}
After expanding the affine abbreviations, all residuals are quadratic.
Let $r_j$ be left side minus right side except that row 8 is negated.
The inherited parent theorem is
\begin{equation}\label{eq:parent}
\exists\mathbf p\in\Pos^{31}\ \bigwedge_{j=1}^{22}r_j(\mathbf p)=0
\quad\longleftrightarrow\quad\neg\Con(\ZFC).
\end{equation}
Negating a residual preserves this theorem. Row 8 now has constant term
$+1$, useful for positivity of the packed polynomial. The all-ones tuple
is not a parent zero: row 5 has residual $1-1(1-1)=1$ there.

The exact parent circuit and a separate sparse description are included.
Our checker proves their equality over the integer polynomial ring,
treating all native ports as independent symbols. This finite identity
does not replace the inherited semantic calibration.

\section{Coefficient extraction using actual supports}\label{sec:packing}
Let $r_1,\ldots,r_m\in\Z[p_1,\ldots,p_n]$ be nonzero polynomials of degree at most two.
Choose positions $w_0=0,w_1,\ldots,w_n>0$ with all unordered pair sums
$w_i+w_j$, including repetitions, distinct. Define
\[
\mathcal S=\{w_i+w_j:0\le i\le j\le n\},\qquad
e(\alpha)=\sum_i\alpha_iw_i.
\]
Write $r_j=\sum_\alpha a_{j,\alpha}p^\alpha$ and
$\mathcal S_j=\{e(\alpha):a_{j,\alpha}\ne0\}$.

\begin{lemma}[Selective target extraction]\label{lem:extraction}
Suppose distinct targets $E_1,\ldots,E_m$ satisfy
\begin{align}
E_j&\ge\max\mathcal S_j,\label{eq:targets1}\\
(E_k-E_j+\mathcal S_j)\cap\mathcal S&=\varnothing\quad(j\ne k).
\label{eq:targets2}
\end{align}
Set
\begin{equation}\label{eq:Dpacking}
\mathcal D(Z)=\sum_{j,\alpha}
(2-|\alpha|)!\prod_i\alpha_i!\;a_{j,\alpha} Z^{E_j-e(\alpha)}.
\end{equation}
Then the exact polynomial identity
\begin{equation}\label{eq:extraction}
[Z^{E_j}]\left(1+\sum_i p_iZ^{w_i}\right)^2\mathcal D(Z)
=2r_j(\mathbf p)
\end{equation}
holds for every $j$. If every coefficient of each $r_j$ has magnitude at
most $H$, every coefficient of $\mathcal D$ has magnitude at most $2H$.
For $p_i\ge0$, every product coefficient has magnitude at most
$2H(1+\sum_i p_i)^2$.
\end{lemma}
\begin{proof}
Pad monomials of degree below two with zero-weight letters. Pair-sum
uniqueness makes the coefficient of $Z^{e(\alpha)}$ in the square
\[
\frac{2!}{(2-|\alpha|)!\prod_i\alpha_i!}\,p^\alpha.
\]
It cancels the multiplier in \eqref{eq:Dpacking} at its own target.
A contribution from row $j$ to a different $E_k$ would require
$E_k-E_j+e(\alpha)\in\mathcal S$, excluded by \eqref{eq:targets2}.

No two terms from different rows have equal reversal exponent, since
$E_j-e(\alpha)=E_k-e(\beta)$ implies the same forbidden membership.
Within a row, uniqueness also prevents collisions. Thus each coefficient
of $\mathcal D$ is one weighted coefficient, with factorial multiplier
at most two. The square has nonnegative coefficients summing to
$(1+\sum p_i)^2$; convolution gives the last bound.
\end{proof}

The full set $\mathcal S$ is indispensable: the packed square includes
pair monomials absent from every residual. The smaller actual support
$\mathcal S_j$ is legitimate only on the source side. Uniform target
spacing greater than $2\max w_i$ is sufficient but unnecessary. We allow
whole support bands to overlap, forbidding contamination only at the
targets actually read.

Use the same 32-element pair-sum set as package 40:
\begin{align*}
(w_0,\ldots,w_{31})={}&(0,10,35,55,62,105,144,145,159,168,232,\\
&245,306,323,325,353,371,374,382,494,499,525,\\
&566,578,610,647,681,685,741,777,783,799).
\end{align*}
The 528 pair sums are distinct. Their finite certificate is independent
of the search that found the set; historically it comes from a
Bose--Chowla construction. With Table~\ref{tab:parent},
\[
t=799,\qquad E_{\max}=7243,\qquad M=E_{\max}+2t=8841.
\]
The independent audit checks all 1974 off-diagonal conditions and
$94\cdot528=49\,632$ dense convolution pairs. It recovers all 22 residuals
over symbolic ports. The coefficient of $Z^{7243}$ in $\mathcal D(Z)$ is exactly 2.

The data order places negated row 8 last; that order is separate from
the original row labels in the table. The target map is explicit.
The new positions were found by allowing insertion before existing
targets. The search has no mathematical authority: the finite conditions
above certify the result.

\section{Translate positivity in the packed word}
A radix digit is nonnegative, while the parent coordinates are positive.
Expanding every substitution $p_i=n_i+1$ in every residual creates
unnecessary records. Instead put
\[
\ell(Z)=\sum_i Z^{w_i},\qquad g(Z)=\sum_i n_iZ^{w_i}.
\]
Then
\[
1+g(Z)+\ell(Z)=1+\sum_i(n_i+1)Z^{w_i},
\]
and Lemma~\ref{lem:extraction} gives
\begin{equation}\label{eq:positivepack}
[Z^{E_j}](1+g(Z)+\ell(Z))^2\mathcal D(Z)
=2r_j(\mathbf n+\mathbf1).
\end{equation}
This represents every positive parent tuple. It does not restrict
coordinates to squares or a sparse subsequence.

There are 94 weighted records, or 99 terms after expanding port
polynomials, instead of 136 records and 161 terms in the shifted baseline.
The register $\ell$ is already required by the low digit mask; inserting
it in the square costs one addition. The finalizer's variable $g$ itself
is unchanged; only the expression inside this square is translated.

Set the fixed bounds
\[
H=8(1+\beta+K_0+M_C+M_F)^2,\qquad\beta_S=2^{2140030}.
\]
Raw coefficients have port degree at most two and coefficient
$\ell^1$ norm at most four, so eight is conservative. Exact integer checks give
\begin{equation}\label{eq:basebound}
\beta_S>4\cdot33^2H,\qquad\beta_S>36H+2.
\end{equation}
The even exponent makes $\beta_S$ a square, a power of two, and 1 modulo
three. The nine nonliteral recipes in $\mathcal D$ and $\beta_S$ are
supplied as actual fixed signed-magnitude binary values, not unspecified indices.

\section{The two-variable source for unrestricted integers}\label{sec:source10}
Letters here through Section~\ref{sec:separator} refer to final coordinates,
not the parent. Start with $f,g\in\Z$ and define polynomial registers
\begin{align}
b&=1+3f,&\mathcal B&=\beta_Sb^2,&\ell&=\sum_i\mathcal B^{w_i},\nonumber\\
N_0&=\mathcal B^{t+1},&N_1&=4\mathcal B^{M+1},&N&=N_0N_1,\label{eq:sourcea}\\
M_c&=N_0-1-(b-1)\ell,&Q_c&=(\mathcal B-2)\sum_j\mathcal B^{E_j},
&T_c&=M_c+Q_cN_0,\nonumber\\
J_c&=2(1+g+\ell)^2\mathcal D(\mathcal B)
+\mathcal B\sum_{i=0}^M\mathcal B^i,
&S_c&=g+J_cN_0,\label{eq:sourceb}\\
R_c&=(S_c+T_c+1)N+T_c+1,&X&=(N-1)R_c,&Y&=N^2,\nonumber\\
\mathscr C&=b\mathcal B^t.&&&&\label{eq:sourcec}
\end{align}
Here $N=4\mathcal B^{9642}$. All sums and powers have fixed indices.

\begin{proposition}[Source interface]\label{prop:source}
\begin{enumerate}[label=(\roman*)]
\item A positive parent zero supplies, for arbitrarily large $f>0$
with $b$ a square and a power of two, an integer $g$ with
$b\le g<\mathscr C$ and $Y\mid\binom{2X}{X}$.
\item If $f\ne0$, $b$ is a square and a power of two,
$0\le g<2\mathscr C$, and $Y\mid\binom{2X}{X}$, the parent has a positive zero.
\item Assuming only $f\ne0$, $b$ square, and $0\le g<2\mathscr C$ gives
$0\le S_c,T_c<N$, $3\mid X$, $X>3b$, and $Y>\max\{b,256\}$.
\end{enumerate}
\end{proposition}
\begin{proof}
First establish broad ranges without assuming binary encoding. A square
$b=1+3f$ with $f\ne0$ has $b\ge4$ and $f>0$.
The leading coefficient of $\mathcal D$ is 2; all others have magnitude
at most $2H$. Since $\mathcal B>2H+1$, $\mathcal D(\mathcal B)>0$.
The radix digits of $M_c$ are $\mathcal B-b$ at selected positions and
$\mathcal B-1$ elsewhere, so $0\le M_c<N_0$. Also
$0\le Q_c<\mathcal B^{7244}<N_1$, whence $0\le T_c<N$.

For the entire interval $0\le g<2b\mathcal B^t$, integrality and
$\ell<2\mathcal B^t$ give $1+g+\ell<3b\mathcal B^t$ and $g<N_0$.
Using \eqref{eq:basebound},
\[
0<J_c<\frac{36Hb^2+\mathcal B}{\mathcal B-1}\mathcal B^{M+1}
<2\mathcal B^{M+1}<N_1.
\]
Thus $0\le S_c<N$. Moreover
$\beta_S\equiv b\equiv\mathcal B\equiv N\equiv1\pmod3$ gives $3\mid X$.
Since $R_c\ge N+1>\mathcal B>b$ and $N\ge4$, $X>3b$.
The fixed radix gives the remaining bound on $Y$.

Now let $b$ be a power of two. The low mask has no binary overlap with
$g$ exactly when $g=\sum_i n_i\mathcal B^{w_i}$ with $0\le n_i<b$.
Put $p_i=n_i+1$. Each coefficient $a_s$ in \eqref{eq:positivepack} obeys
\[
|a_s|\le2H(1+\textstyle\sum_i p_i)^2<2H(33b)^2<\mathcal B/2.
\]
Thus the radix digits of
\[
K_c=(1+g+\ell)^2\mathcal D(\mathcal B)
+(\mathcal B/2)\sum_{s=0}^M\mathcal B^s
\]
are exactly $\mathcal B/2+a_s$, strictly between zero and $\mathcal B$.
This is a proof abbreviation; the circuit evaluates $J_c=2K_c$ without
division. At target $E_j$ the digit is $\mathcal B/2+2r_j(\mathbf p)$.
For $0<d<\mathcal B$, no overlap with $\mathcal B/2-1$ holds precisely
when $d=\mathcal B/2$. The high mask $Q_c$, after a common shift by one
binary bit, applies this test at every target. Low and high blocks are
disjoint, so
\[
\text{no carries in }S_c+T_c
\ \longleftrightarrow\
g\text{ encodes digits and }r_j(\mathbf n+\mathbf1)=0\text{ for all }j.
\]
Lemma~\ref{lem:binomial} below identifies this with
$N^2\mid\binom{2X}{X}$, proving (ii).

Conversely, given a parent zero, set $n_i=p_i-1$ and choose arbitrarily
large $b=2^{2q}>\max_i n_i$, so $f=(b-1)/3$ is positive.
The all-ones tuple is not a parent zero, so some $n_i>0$ and
$g=\sum n_i\mathcal B^{w_i}\ge\mathcal B\ge b$.
A geometric-series bound gives $g<b\mathcal B^t$. Both masks pass,
giving (i).
\end{proof}

\begin{lemma}[Central binomial carry identity]\label{lem:binomial}
For $N=2^e$, $0\le S,T<N$, and
$X=(N-1)((S+T+1)N+T+1)$,
\[
N^2\mid\binom{2X}{X}
\quad\longleftrightarrow\quad
\text{binary addition of }S,T\text{ has no carries}.
\]
\end{lemma}
\begin{proof}
The exact base-$N$ expression is
$X=(S+T)N^2+(N-1-S)N+(N-1-T)$.
Hence $\pc(X)=2e+\pc(S+T)-\pc(S)-\pc(T)=2e-\tau_2(S,T)$,
where $\tau_2$ counts carries with multiplicity. The factorial valuation
formula gives $\nu_2\binom{2X}{X}=\pc(X)$, proving the assertion.
\end{proof}

\section{The explicit ten-integer polynomial}\label{sec:ten}
The ten unknowns are
\[
(f,g,h,k,l,w,x,y,z,m)\in\Z^{10}.
\]
Use the source quantities $b,\mathcal B,\ell,N,X,Y,\mathrm{cap}$
from Section~5, writing $\mathrm{cap}:=\mathscr C$. Letters introduced below are polynomial abbreviations,
not additional unknowns. In particular, $g$ remains the original unknown
in every expression below; the successor translation is confined to
the source expression $(1+g+\ell)^2\mathcal D(\mathcal B)$.

Define the following registers:
\begin{align*}
 L&=lY,& U&=2LX,& V&=4gwY,& W&=bw,\\
 A&=U(V+1),& B_1&=2X+1,& C&=B_1+(A-2)h,\\
 P&=U^2V,& K&=X+1+k(P-2),& S&=2A-5,\\
 T&=3WC-2(W^2-1),&
 r_1&=(P^2-4)K^2+4.
\end{align*}
The range margin and positivity register are
\begin{align*}
 \Delta_b&=8\,\mathrm{cap}^3gK^2
 -g^2\bigl(32(C-KL)^2\mathrm{cap}^3+g^2K^2\bigr),\\
 O&=f^2l^2\Delta_b,\qquad \xi=zx.
\end{align*}
The modified Lucas certificate is
\begin{align*}
 D&=(A^2-4)C^2+4,& E&=C^2D\xi,\\
 F&=4(A^2-4)E^2+1,&
 G_*&=2(A+2)(A-2)^2E^2-1,\\
 H_*&=C+(B_1-C)F+16EFy,&
 I_*&=(G_*^2-1)H_*^2+1,\\
 r_2&=DFI_*.
\end{align*}
Set
\begin{align*}
 r_0&=b\bigl(1+(3O-4)b^2z^2\bigr),\\
 \Sigma&=\xi S^5,\qquad r=T+Sy,\\
 \Theta&=(r^3-rS^2-S^3)(r^2+23S^2),\\
 Z_{\rm n}&=\Sigma m-\Theta,\\
 a_0&=\Sigma^2r_0,\qquad
 a_1=\Sigma^2b^2r_1,\qquad
 a_2=\Sigma^2b^2K^2r_2.
\end{align*}
Finally put
\begin{equation}\label{eq:ten-polynomial}
\boxed{\begin{aligned}
 V_{\rm n}&=Z_{\rm n}^2+a_2-a_0-a_1,\\
 W_{\rm n}&=V_{\rm n}^2+4(Z_{\rm n}^2a_2-a_0a_1),\\
 F_{10}&=W_{\rm n}^2-16V_{\rm n}^2Z_{\rm n}^2a_2.
\end{aligned}}
\end{equation}
The requested unrestricted-integer equation is $F_{10}=0$.
The complete expanded circuit, including fixed-power schedules, is
\file{data/polynomial.json}. Thus this display specifies a single
polynomial, despite its readable sequence of abbreviations.

\subsection{The norm identity}
For indeterminates $t,a_0,a_1,a_2$, let
\[
 \mathcal N(t,a_0,a_1,a_2)
 =\bigl((t^2+a_2-a_0-a_1)^2+4(t^2a_2-a_0a_1)\bigr)^2
 -16(t^2+a_2-a_0-a_1)^2t^2a_2.
\]
Multiplication of conjugate factors proves the exact identity
\begin{equation}\label{eq:norm}
\mathcal N(t,a_0,a_1,a_2)
 =\prod_{\varepsilon\in\{-1,1\}^3}
 \bigl(t+\varepsilon_0\sqrt{a_0}
         +\varepsilon_1\sqrt{a_1}
         +\varepsilon_2\sqrt{a_2}\bigr).
\end{equation}
It is valid over the complex numbers, including zero radicands and
coincident squareclasses. All radicals cancel; the result is an integer
polynomial monic of degree eight in $t$. The audit separately expands
both sides to the same 35-term polynomial.

The conceptual weights in this construction are $(1,b,bK)$.
The previous package-40 weights were
$(1,r_0,r_0(P^2+K^2+3))$. The smaller weights are valid because of
arithmetic information about these particular radicands. They are
not a universal square-test identity for arbitrary inputs.

\section{An arithmetic square-root separator}\label{sec:separator}
\begin{lemma}[Coprimality separates one squareclass]\label{lem:coprime}
Let $b,J_0,A_0,A_1,K\in\Z$, with $|b|>1$ and
$\gcd(b,J_0)=1$. If a signed sum
\[
 \sqrt{bJ_0}\ \mathbin{\pm}\ b\sqrt{A_0}\
                  \mathbin{\pm}\ bK\sqrt{A_1}
\]
is rational, then $bJ_0$ is a nonnegative integer square.
\end{lemma}
\begin{proof}
If $bJ_0$ is not such a square, write $bJ_0=da^2$ with $a>0$ and
$d\ne1$ a signed squarefree integer. Distinct signed squarefree radical
classes are linearly independent over $\mathbb Q$: they are distinct
characters in the multiquadratic extension generated by the radicals.
The coefficient of $\sqrt d$ must therefore cancel. Every contribution
from the last two terms is an integer multiple of $b$, so $b\mid a$.
Then $b^2\mid bJ_0$, hence $b\mid J_0$, contradicting coprimality.
This argument also covers negative $d$.
\end{proof}

\begin{lemma}[A Pell remainder prevents two-root cancellation]\label{lem:remainder}
Suppose $P,K,R\in\Z$, $K\ne0$, and $4\mid R$. Put
$A_0=(P^2-4)K^2+4$. If a signed sum
$\sqrt{A_0}\mathbin{\pm}K\sqrt R$ is rational, then both $A_0$ and $R$
are nonnegative integer squares.
\end{lemma}
\begin{proof}
Call the signed rational sum $q$. If $q\ne0$, squaring gives a rational
expression for $\sqrt{A_0}$ with denominator $2q$, and then one for
$\sqrt R$. A rational square root of an integer is an integer.
If $q=0$, then $A_0=K^2R$, so $K^2\mid4$. For $|K|=1$ this gives
$A_0=R=P^2$. For $|K|=2$ it gives $R=P^2-3$, which is $1$ or $2$
modulo $4$, contradicting $4\mid R$.
\end{proof}

\begin{proposition}[All square predicates are recovered]\label{prop:separation}
Every integer zero of $F_{10}$ satisfies
\[
 \Sigma\ne0,\quad \Sigma\mid\Theta,\quad S\mid T,\quad
 b\text{ is a positive square},\quad O>0,\quad
 r_1,r_2\text{ are nonnegative squares}.
\]
\end{proposition}
\begin{proof}
The integer $S=2A-5$ is odd and nonzero. The polynomial
\[
 \pi(v)=(v^3-v-1)(v^2+23)
\]
has no rational root. Consequently
$\Theta=S^5\pi((T+Sy)/S)\ne0$ for every integer assignment.
At $\Sigma=0$ equation~\eqref{eq:ten-polynomial} is $\Theta^8$.
Thus a zero has $\Sigma\ne0$, whence $\xi,z,x\ne0$.

Divide the norm equation by $\Sigma^8$.
The rational number $m-\Theta/\Sigma$ is a root of a monic integer
polynomial with radicands $r_0,b^2r_1,b^2K^2r_2$. The rational-root
theorem gives $\Sigma\mid\Theta$. The product identity also gives a
rational signed sum
\begin{equation}\label{eq:rational-sum}
 \varepsilon_0\sqrt{r_0}
 +\varepsilon_1b\sqrt{r_1}
 +\varepsilon_2bK\sqrt{r_2}.
\end{equation}
Signs of integer weights may be absorbed into the independent signs.

Write $J_0=1+(3O-4)b^2z^2$. Then $\gcd(b,J_0)=1$.
Since $b=1+3f$, it is nonzero and is never $-1$. If $|b|>1$,
Lemma~\ref{lem:coprime} implies that $r_0=bJ_0$ is a square.

We handle the remaining possibility $b=1$ before assuming positivity.
It gives $f=O=0$ and $r_0=1-4z^2$, an odd negative integer.
Both $r_1$ and $r_2$ are divisible by four: $P,A$ are even and
$D=(A^2-4)C^2+4$. Write $r_0=da^2$ with signed squarefree $d\ne1$.
Both $d$ and $a$ are odd. A radical from either of the last two terms
of~\eqref{eq:rational-sum} in this same squareclass has an even integer
coefficient. Its cancellation would require an odd integer plus even
integers to vanish. Thus $b=1$ is impossible.

Also $r_0\ne0$. Indeed, $J_0=0$ would imply
$(3O-4)b^2z^2=-1$, hence $|b|=|z|=1$ and $O=1$.
But then $b=1,f=0,O=0$, a contradiction.
Thus $r_0$ is a positive square. Unique factorization and coprimality
show that $b,J_0$ are signed squares with the same sign.
A negative nonzero square is never $1$ modulo $3$, so $b,J_0>0$.
If $O\le0$, then $J_0\le1-4(bz)^2<0$, a contradiction.
Therefore $O>0$, $f,l\ne0$, and $\Delta_b>0$.

Since $\mathrm{cap}=b\mathcal B^t>0$, the case $K=0$ would give
$\Delta_b=-32g^2C^2\mathrm{cap}^3\le0$. Thus $K\ne0$.
Subtract the integral first root in~\eqref{eq:rational-sum} and divide
by $b$. Lemma~\ref{lem:remainder}, with $R=r_2$, proves that $r_1,r_2$
are squares. No source size bound, preliminary positivity of $g$,
or Lucas recovery was used in these steps.

It remains to justify $S\mid T$. If not, choose a prime with
$\alpha=\nu_p(S)>\beta=\nu_p(T)$. In the two homogeneous factors
of $\Theta$, the terms $r^3$ and $r^2$ have uniquely smallest valuations,
respectively $3\beta$ and $2\beta$, because
$\nu_p(r)=\nu_p(T+Sy)=\beta$.
It follows that $\nu_p(\Theta)=5\beta<5\alpha$,
contradicting $S^5\mid\Theta$.
\end{proof}

\subsection{The inherited Lucas bridge and its precise use}
The following are the existing Sun/package-31 bridge statements, used
in the parametric form justified in package 40 \cite{repo40,sun2026}.
Let $u_0(A)=0,u_1(A)=1$, and
$u_{n+1}(A)=Au_n(A)-u_{n-1}(A)$.
For a source satisfying Section~5:
\begin{enumerate}
\item A positive parent zero supplies positive bridge witnesses with
$A>2$ even, $B_1>0$ odd,
\[
 b\text{ square},\quad b\le g<\mathrm{cap},\quad
 C=u_{B_1}(A),\quad r_1\text{ square},\quad S\mid T,\quad O>0.
\]
\item For arbitrary signed inputs, the conditions
\[
 b\text{ square},\quad O>0,\quad \xi\ne0,\quad
 r_1,r_2\text{ squares},\quad S\mid T
\]
imply that $b$ is a power of two, $0<g<2\,\mathrm{cap}$, and
$Y\mid\binom{2X}{X}$.
\item For $A>2$ even, $B_1>0$ odd, and $C=u_{B_1}(A)$, every positive
integer $d$ admits $\xi>0$ divisible by $d$ such that each residue class
modulo $\xi$ contains a $y$ making $r_2$ a square.
\end{enumerate}
These are inherited arithmetic theorems, not outputs of a finite audit.
The full original proof and source-dependent hypotheses are included in
\file{upstream/PACKAGE40_REPORT.tex} and
\file{upstream/SIGNED10_PROOF.md}.

The new source preserves every premise used in those parametric proofs:
$b=1+3f$, the forward and reverse central-binomial conditions, the
pre-decoding bounds $3\mid X$, $X>3b$, and $Y>\max(b,256)$, and
\[
 \mathrm{cap}=2^{1709883970}b^{1599}.
\]
The fixed exponent of two is even and the exponent of $b$ is odd,
as required by the inherited positive bridge. The new target schedule
changes $X,Y$ and their degrees, but not this cap.

\begin{theorem}[Ten-integer equivalence]\label{thm:ten}
The polynomial of~\eqref{eq:ten-polynomial} has a zero in $\Z^{10}$
if and only if the positive quadratic parent has a zero.
Consequently its nonsolvability is provably equivalent in ZFC to
$\Con(\ZFC)$.
\end{theorem}
\begin{proof}
For soundness, Proposition~\ref{prop:separation} gives all hypotheses
of the reverse bridge. Source soundness then decodes a positive parent
zero. For clarity, $O>0$ itself implies $g>0$ and
\[
 0<g<2\,\mathrm{cap},\qquad 4(C-KL)^2<K^2.
\]
Indeed $g\le0$ makes $\Delta_b\le0$; with $g>0$ its positivity yields
$g^3<8\,\mathrm{cap}^3$ and
$4g(C-KL)^2<K^2$. This supplies the broad source range before binary
base recovery. There is no assumption that $g$ already encodes digits.

Conversely, the forward bridge supplies the witnesses in item 1.
The positive integer $\Delta_b$ is a multiple of $g$, so
$O=f^2l^2\Delta_b\ge g\ge b\ge4$.
Then $d=3O-4>0$ is $2$ modulo $3$, hence not a square.
A positive Pell unit for $d$, raised to a suitable finite order
modulo $b$, gives $u^2-dv^2=1$ with $v>0$ and $b\mid v$.
Taking $z=v/b>0$ makes $r_0=b u^2$ a square.

The polynomial $\pi$ has a root modulo every positive integer.
For completeness, its cubic factor has discriminant $-23$.
At an odd prime other than $23$, either the cubic has a simple root,
or its irreducibility makes the discriminant a square and the
quadratic has a simple root. At $23$, the cubic has the simple root $3$.
At $2$, $-23\equiv1\pmod8$ is a $2$-adic square. Hensel lifting and
the Chinese remainder theorem give the assertion.

Now choose $\xi>0$ divisible by $z$ using item 3, and put $x=\xi/z$.
Write $T=Sq$. Choose a root $c$ of $\pi$ modulo $\xi$ and prescribe
$y\equiv c-q\pmod\xi$. Interpolation supplies such a $y$ with $r_2$
square, while
$\Theta=S^5\pi(q+y)$ is divisible by $\Sigma=\xi S^5$.
All three weighted square roots are integers. Choose $m$ so that
one factor of~\eqref{eq:norm} vanishes. This proves completeness.
Finally apply the inherited parent equivalence~\eqref{eq:parent}.
\end{proof}

\subsection{Exact cost and degree}
The delivered circuit contains 189 additions, 25 subtractions, and
495 multiplications, totaling 709 live binary operations.
It shares $\Sigma^2b^2$, $K^2$, and $Z_{\rm n}^2a_2$.
The displayed factorization of $W_{\rm n}$ saves a scalar multiplication.
The conceptual product $bK$ need not be computed as a separate gate.

Degree propagation gives the total-degree upper bound
$54\,768\,104$.
Specialize all ten variables to one indeterminate $T$ and compute
the leading coefficient at that upper bound using the actual fixed
coefficient tuple. Modulo $1\,000\,000\,007$ it is
\[
 323\,596\,601\ne0.
\]
This proves that the bound is attained over $\Z$, hence that it is the
exact total degree. This is exact coefficient arithmetic, not a
probabilistic identity test. Independent literal-formula replay also
checks nonzero leaders modulo two further primes.

\section{The nine-positive-variable companion}\label{sec:nine}
This branch combines the same new quadratic source with the existing
Jones ratio finalizer in package 45 \cite{repo45}. Its raw positive
parent is the package-40 parent in Table~\ref{tab:parent}; package 45
used an affine-shifted parent. The source interface, rather than a claim
of pointwise equality between those parents, justifies the composition.

\subsection{Source and domain}
Initially take mixed coordinates
\[
 \epsilon,\phi,n\in\N,\qquad g,h,s,w,i,j\in\Pos.
\]
Define the fixed constants
\[
 M_{\rm coef}=2^{2134927},\qquad b_{\rm floor}=2^{2134939}
 =4096M_{\rm coef},
\]
and $b=b_{\rm floor}+\epsilon$, $\mathcal B=b^3$.
Every weighted coefficient of the same $\mathcal D$ has absolute value
less than $M_{\rm coef}$; the exact recipes and bit-length audit certify
this bound. With the same $t=799$, $E_{\max}=7243$, and $M=8841$, set
\begin{align*}
 \ell&=\sum_{i=1}^{31}\mathcal B^{w_i},&
 N_0&=\mathcal B^{t+1},&
 N_1&=4\mathcal B^{M+1},&
 N&=N_0N_1,\\
 M_c&=N_0-1-(b-1)\ell,&
 Q_c&=(\mathcal B-2)\sum_{j=1}^{22}\mathcal B^{E_j},\\
 T_c&=M_c+Q_cN_0,\\
 J_c&=2(1+g+\ell)^2\mathcal D(\mathcal B)
       +\mathcal B\sum_{a=0}^{M}\mathcal B^a,&
 S_c&=g+J_cN_0,\\
 R&=(N-1)\bigl((S_c+T_c+1)N+T_c+1\bigr),&
 G_{\rm cap}&=32b\mathcal B^t.
\end{align*}

These formulas have unconditional positivity on the mixed domain.
In particular $\mathcal D(\mathcal B)>0$ because its leading coefficient
is $2$ and $\mathcal B-1>M_{\rm coef}$.
Hence $R,N\ge8$, $b\le R,N$, and $b,G_{\rm cap}>0$.

The broad range needed for soundness is $0<g<4b\mathcal B^t$.
Integrality and $\ell<2\mathcal B^t$ give
\[
 1+g+\ell<5b\mathcal B^t.
\]
The sparse contribution to $J_c$ is less than
\[
 \frac{50M_{\rm coef}b^2\mathcal B^{M+1}}{\mathcal B-1}.
\]
The bound $b\ge4096M_{\rm coef}$ gives
$50M_{\rm coef}b^2<\mathcal B/64$.
Adding the geometric mask therefore leaves $J_c<N_1$,
and $0<S_c,T_c<N$.
This is the only changed numerical estimate in the inherited source
argument; using the old numerator $32$ after adding $\ell$ would not
justify the broad-range bound.

If $b$ is a power of two, the no-carry proof applies.
It recovers digits $0\le z_i<b$ and positive witnesses $p_i=z_i+1$.
The coefficients to be centered are bounded by
\[
 M_{\rm coef}\left(1+\sum_i p_i\right)^2
 \le M_{\rm coef}(32b)^2
 \le\mathcal B/4<\mathcal B/2.
\]
The independent target certificate supplies exact isolation, so the
high mask forces all $r_j(\mathbf p)=0$.
Conversely a positive parent zero permits an arbitrarily large binary
$b\ge b_{\rm floor}$ and
$g=\sum_i(p_i-1)\mathcal B^{w_i}>0$ with
$g<b\mathcal B^t$. Thus the exact interface is
\begin{align*}
 \text{parent zero}&\ \Longrightarrow\
 \text{binary }b,\quad N^2\mid\binom{2R}{R},\quad32g<G_{\rm cap},\\
 \text{binary }b,\quad N^2\mid\binom{2R}{R},\quad8g<G_{\rm cap}
 &\ \Longrightarrow\ \text{parent zero}.
\end{align*}

\subsection{The full finalizer}
All notation in this subsection is local to the nine-variable branch.
Define
\begin{align*}
 U&=N^2w,&Y&=N^2s,&M_0&=RY,&P&=2M_0^2U,\\
 K&=R+1+h(P-1),&A&=M_0(U+1),&
 B_1&=2R+1,&W_1&=bw,\\
 C&=B_1+W_1+\phi,&J_*&=4A-5,&
 Q&=3W_1C-2(W_1^2-1),\\
 D&=(A^2-1)C^2+1,&E&=iJ_*DC^2,&
 F&=(A^2-1)E^2+1,\\
 G&=A+F(CD+1-A),&H&=B_1+jC,&
 I&=(G^2-1)H^2+1,\\
 a_0&=(P^2-1)K^2+1,&a_1&=DFI,&
 u&=J_*F,&v&=QF+J_*(H-C).
\end{align*}
The ratio margin and strict square-root majorants are
\begin{align*}
 \mu&=\bigl(K^2-4(C-KY)^2\bigr)G_{\rm cap}-8gK^2,\\
 V_0&=PK,\qquad V_1=(AC)(AE)(GH),\qquad M_{\rm div}=Q+H,\\
 \omega&=u\mu,\qquad
 \theta=u\bigl(n+M_{\rm div}-\mu(2M_{\rm div}+V_0V_1)\bigr)+v,\\
 b_0&=\omega^2a_0,\qquad b_1=\omega^2V_0^2a_1.
\end{align*}
The equation is
\begin{equation}\label{eq:nine-polynomial}
 \boxed{F_9=(\theta^2-b_0-b_1)^2-4b_0b_1=0.}
\end{equation}

To obtain the strictly positive version, replace the mixed coordinates
$\epsilon,\phi,n$ by their respective positive coordinates minus one.
To obtain the nonnegative version, retain those three coordinates and
replace each of $g,h,s,w,i,j$ by its nonnegative coordinate plus one.
These are domain bijections. The equation is not asserted to have the
same consistency property when its inputs are allowed to be arbitrary
signed integers.

\subsection{Why one equation enforces the ratio conditions}
The inherited Jones ratio theorem has the following form.
Under $R,N\ge8$, $b\le R,N$, and $b,G_{\rm cap}>0$, the conditions
\[
 a_0,a_1\text{ squares},\qquad u\mid v,\qquad\mu>0
\]
imply binary $b$, $N^2\mid\binom{2R}{R}$, and $8g<G_{\rm cap}$.
Conversely, binary $b$, the divisibility, and $32g<G_{\rm cap}$
supply positive $h,s,w,i,j$ and nonnegative $\phi$ with those conditions.
The complete dependency identification and inherited finalizer are in
\file{upstream/PACKAGE45_NINE_PROOF.md}.
There are also existing Lean proof sources for
\file{Jones1982.Nine.ratio_sound} and
\file{Jones1982.Nine.ratio_complete} in
\file{Diophantine.Paper1982.NineRatio}.
Their signatures have arithmetic hypotheses, with no machine or compiler
assumption. The archive includes their 20-file local import closure and
a source-hash audit under \file{upstream/JONES_RATIO_LEAN/}.
This is inspected existing proof source; a new Lean kernel replay was
not performed, and Mathlib remains an external dependency.

For completeness we give the algebraic finalizer proof, since it
distinguishes a valid domain translation from merely adding variables.
On every mixed-domain assignment, all of $P,K,A,C,J_*,D,E,F,G,H$ exceed
one, while $a_0,a_1,u>0$. The exact identities
\begin{align*}
 V_0^2-a_0&=K^2-1>0,\\
 (AC)^2-D&=C^2-1>0,\\
 (AE)^2-F&=E^2-1>0,\\
 (GH)^2-I&=H^2-1>0
\end{align*}
give $0<a_0<V_0^2$ and $0<a_1<V_1^2$.
Moreover $Q>0$, $H-C>0$, and
\[
 0<\frac vu=\frac Q{J_*}+\frac{H-C}{F}<M_{\rm div}-1.
\]
If $\mu=0$, then $\omega=0$ and
$\theta/u=n+M_{\rm div}+v/u>0$, excluding a zero of~\eqref{eq:nine-polynomial}.

Otherwise the four-factor norm identity gives
$\theta/\omega=r_0+V_0r_1$ with $r_a^2=a_a$.
The roots are real. The rational sum cannot vanish, since that would
give $|r_1|=|r_0|/V_0<1$, contradicting the positive integer $a_1$.
Squaring the nonzero rational sum makes both roots rational and hence
integers. The strict bounds yield $|r_a|\le V_a-1$, so
\[
 Z=2M_{\rm div}+V_0V_1+r_0+V_0r_1\ge2M_{\rm div}+1>0.
\]
The vanishing factor now says
\[
 n+M_{\rm div}+v/u=\mu Z.
\]
The left side is positive, giving $\mu>0$.
Its remaining terms are integers, giving $u\mid v$.
Conversely, from the two squares, $u\mid v$, and $\mu>0$, choose integer
roots and set $n=\mu Z-M_{\rm div}-v/u$. The same bounds give $n\ge2$,
and a norm factor vanishes.

The source interface and the ratio theorem now prove that
\eqref{eq:nine-polynomial} has a zero exactly when the positive parent
does. This proves the nine-variable cases of Theorem~\ref{thm:main}.

\subsection{Exact degree despite cancellation}
Put $d_N=3(E_{\max}+3t+2)=28\,926$.
On the diagonal the source has
$N$ of degree $d_N$ and leading coefficient $4$, and
$R$ of degree $3d_N$ and leading coefficient $16$.
The new sparse translation is strictly lower degree than the geometric
mask, so the inherited finalizer's degree ledger still applies.

Naive propagation in~\eqref{eq:nine-polynomial} would conceal a leading
cancellation. Resolve it with the exact identity
\begin{align*}
 V_1^2-a_1={}&(C^2-1)(AE)^2(GH)^2\\
 &+D(E^2-1)(GH)^2+DF(H^2-1).
\end{align*}
The first two summands each have degree $298d_N+64$ and diagonal
leading coefficient $2^{518}$; the third has lower degree.
With
$E_0=u(n+M_{\rm div})+v-2\omega M_{\rm div}$, exactly
\[
 \theta^2-\omega^2(a_0+V_0^2a_1)
 =E_0^2-2E_0\omega V_0V_1
  +\omega^2V_0^2(V_1^2-a_1)-\omega^2a_0.
\]
The third term uniquely dominates, with degree
$576d_N+6t+140$ and leading coefficient $2^{1025}$.
The final subtracted term in~\eqref{eq:nine-polynomial} has lower degree.
Therefore
\[
 \deg F_9=1152d_N+12t+280=33\,332\,620,
 \qquad [T^{33\,332\,620}]F_9(T,\ldots,T)=2^{2050}.
\]
The coordinate shifts preserve this leading term.
The primary positive evaluator uses 682 operations; its natural
companion uses 685. Separate 700/703-operation circuits make the
degree cancellation explicit. Their greater costs are not substituted
for the primary evaluator counts.

\section{A squared-modulus quotient: a 66-operation equation}
\label{sec:short}

The constructions with fewest variables need large arithmetic circuits.
There is a separate improvement at the short-circuit end of the comparison:
the inherited fifteen-positive-variable equation can be evaluated in
\textbf{66 binary arithmetic operations}.  Its exact degree is 157.  This
saves one operation against the repository's 67-operation, degree-155
baseline.  The same idea lowers the costs of the degree-80 and degree-68
alternatives from 73 and 75 to 72 and 74 operations, respectively.
The saving comes from a stronger auxiliary congruence, followed by a change
of the quantified quotient.  It is not an identity between the old and new
polynomials at the same quotient coordinate.

\subsection{The explicit polynomial and its fixed coefficients}

Use the entire coefficient tuple of package 36 at the source commit
\texttt{124c07c0b7}: the ports are
\(\beta=B-1,K_0,M_C,M_F\), with \(B=2^{1012115}\).
Here \(K_0\) is the tuple's \texttt{K} and \(M_F\) is its already shifted
\texttt{MF\_source}.  The four ports are fixed integers, not additional
unknowns.  The coefficient construction and its machine semantics are
inherited; the argument below introduces no rebinding or change of masks.

All fifteen unknowns are strictly positive integers:
\[
 (J,F,\alpha,z,f,h,i,T,s,w,\tau,\eta,\zeta,y,\gamma)\in\Pos^{15}.
\]
Define the following polynomial abbreviations, which introduce no new
unknowns:
\begin{align*}
 q&=\beta J+1,                 & C&=q-F-\alpha,\\
 X&=wq,                       & Y&=sq^3, & E&=XY,\\
 k&=\eta+\zeta,               & c&=kY+\eta, & a&=E+Y,\\
 H&=4a+3,                     & \Delta&=a^2+H, & A&=a+2,\\
 D&=X+ac+\gamma H,            & G&=(q-1)(q-F)+\alpha,\\
 R&=G(q^2-1)+(M_C+qM_F)J,     & S&=i\Delta c^2,\\
 V&=f^2(cT-R)-c.&&&
\end{align*}
In particular, \(\Delta=A^2-1=(a+1)(a+3)>0\) on the positive domain.
Set
\begin{align*}
 N_0&=\tau^2-(XY^2)(XY^2+1)k^2,\\
 N_1&=D^2-\Delta c^2,\\
 N_a&=S^2(V^2-y^2)+y^2,\\
 N_k&=k-hE-R,\\
 N_t&=(K_0+w)C+(q-F)-z(q-1),\\
 N_{\mathrm{str}}&=\Delta f^2-S^2.
\end{align*}
The proposed equation is the ordinary integer polynomial equation
\begin{equation}
 \boxed{\quad
 P_{66}:=N_0N_1N_aN_kN_tN_{\mathrm{str}}-\Delta=0.
 \quad}
 \label{eq:p66}
\end{equation}
Every exponent in this definition is fixed.  The divisions and Pell
polynomials used below occur only in the proof and witness construction.

\subsection{Why a square modulus permits a cheaper quotient}

Write \(\chi_t(n),\psi_t(n)\) for the integer Pell polynomials determined by
\[
 \chi_t(0)=1,\qquad \psi_t(0)=0,
\]
\[
 \chi_t(n+1)=t\chi_t(n)+(t^2-1)\psi_t(n),\qquad
 \psi_t(n+1)=\chi_t(n)+t\psi_t(n).
\]
Thus \(\chi_t(n)^2-(t^2-1)\psi_t(n)^2=1\).

\begin{lemma}[Odd Pell quotient identities]
\label{lem:odd-quotient}
For each \(r\in\N\) there is a polynomial \(Q_r\in\Z[Z]\) such that
\begin{align}
 \chi_t(2r+1)&=tQ_r(t^2),\label{eq:odd-quotient-identity}\\
 Q_r(0)&=(-1)^r(2r+1),\label{eq:odd-quotient-constant}\\
 Q_r(1-A^2)&=(-1)^r\psi_A(2r+1).
 \label{eq:odd-quotient-conjugate}
\end{align}
\end{lemma}

\begin{proof}
The odd-index subsequence gives
\[
 Q_0(Z)=1,\qquad Q_1(Z)=4Z-3,\qquad
 Q_{r+1}(Z)=(4Z-2)Q_r(Z)-Q_{r-1}(Z).
\]
This proves \eqref{eq:odd-quotient-identity} and
\eqref{eq:odd-quotient-constant} by induction.  The sequence
\(\psi_A(2r+1)\) starts with \(1,4A^2-1\) and has recurrence coefficient
\(4A^2-2\).  Substituting \(Z=1-A^2\) changes the coefficient of the
displayed recurrence to its negative, giving
\eqref{eq:odd-quotient-conjugate} by the same induction.  These are
identities over the integers and therefore also hold modulo any positive
integer.
\end{proof}

\begin{lemma}[Squared-modulus reconstruction]
\label{lem:squared-modulus-reconstruction}
Suppose \(A>1\), \(R>0\), \(R\equiv3\pmod4\),
\(c=\psi_A(R)>0\), and \(i,f>0\) satisfy
\[
 \Delta=A^2-1,\qquad f^2=1+\Delta i^2c^4.
\]
Put \(S=\Delta ic^2\).  There are positive integers \(V,y,T\) satisfying
\[
 S^2(V^2-y^2)+y^2=1,\qquad V=f^2(cT-R)-c.
\]
\end{lemma}

\begin{proof}
Take
\[
 V=\frac{\chi_S(R)}{S},\qquad y=\psi_S(R).
\]
The quotient is integral by \eqref{eq:odd-quotient-identity}; both values
are positive because \(S>1\) and \(R>0\).  The Pell identity gives the
required norm equation.  Also
\[
 S^2=\Delta(f^2-1)\equiv1-A^2\pmod{f^2}.
\]
Since \((R-1)/2\) is odd, Lemma~\ref{lem:odd-quotient} gives the two
congruences
\begin{equation}
 V\equiv-c\pmod{f^2},\qquad V\equiv-R\pmod c.
 \label{eq:squared-modulus-congruences}
\end{equation}
For the second, use \(c\mid S\) and
\eqref{eq:odd-quotient-constant}.  The strong equation gives
\(f^2\equiv1\pmod c\), hence \(\gcd(c,f)=1\).
The numerator \(V+c+Rf^2\) is therefore divisible by both \(f^2\) and
\(c\), and thus by \(cf^2\).  Consequently
\[
 T=\frac{V+c+Rf^2}{cf^2}
\]
is a positive integer.  Rearranging proves the stated quotient identity.
\end{proof}

The earlier auxiliary construction used congruence modulo \(f\).  Retaining
the full modulus \(f^2\) is what makes the old quotient divisible by
\(f\), after choosing a suitable auxiliary Pell ordinate.

\subsection{Soundness, completeness, and the native congruence}

\begin{theorem}[Existence-equivalence of the short equation]
\label{thm:p66-consistency}
With the authentic package-36 coefficient tuple, the inherited ordinary
mathematical compiler and K215 calibration imply
\[
 \mathrm{ZFC}\vdash
 \left[
 \neg\exists\mathbf{x}\in\Pos^{15}\;P_{66}(\mathbf{x})=0
 \right]
 \;\longleftrightarrow\;
 \operatorname{Con}(\mathrm{ZFC}).
\]
The additional reduction from the inherited explicit polynomial to
\eqref{eq:p66} requires no consistency assumption.
\end{theorem}

\begin{proof}
Let \(P_{67}\) denote the inherited polynomial on the same fifteen slots.
Its only different abbreviation is
\[
 V_{\mathrm{old}}=c(T_{\mathrm{old}}f-1)-Rf^2.
\]
The exact all-value identity
\begin{equation}
 P_{66}(\ldots,f,T,\ldots)
   =P_{67}(\ldots,f,fT,\ldots)
 \label{eq:p66-pullback}
\end{equation}
follows by substitution.  Thus every positive new zero gives a positive old
zero, since \(fT>0\).  Inherited soundness supplies a finite ZFC refutation.

Conversely, a finite refutation supplies a positive old zero.  Its inherited
Pell and compiler recovery gives
\[
 c=\psi_A(R),\qquad f^2=1+\Delta i^2c^4,
 \qquad q=B^N,\qquad J=1+B+\cdots+B^{N-1}
\]
for some positive integer \(N\), together with the native bit condition
\[
 (C-1)\mathbin{\mathrm{AND}}(M_CJ+1)=0.
\]
The tuple has \(4\mid B\) and \(M_C\equiv2\pmod4\).  Hence
\(q\equiv0\pmod4\), \(J\equiv1\pmod4\), and
\(M_CJ+1\equiv3\pmod4\).  The two low bits of the nonnegative integer
\(C-1\) must vanish, so \(C\equiv1\pmod4\).  Using the exact identity
\(G=q(q-F)-C\), the definition of \(R\) yields
\[
 R\equiv C+2\equiv3\pmod4.
\]
This phase follows from the authentic native decoding; it is not a
restriction silently imposed on arbitrary coefficient assignments.

Apply Lemma~\ref{lem:squared-modulus-reconstruction} at these same values
of \(A,R,c,i,f\).  Replace only the supplied auxiliary ordinate \(y\) and
quotient \(T\) by the lemma's positive integers.  The auxiliary norm factor
is again one.  The other factors are unchanged; in particular the inherited
unit classification gives \(N_0=N_1=N_k=N_t=1\) and
\(N_{\mathrm{str}}=\Delta\).  Thus \(P_{66}=0\).

We have established equivalence of positive solvability, which composes
with the inherited refutation calibration.  Negating gives the displayed
consistency equivalence.  All new steps are finite integer algebra,
induction on polynomial recurrences, and elementary divisibility, and can
be carried out in ZFC.
\end{proof}

The reverse map may change both auxiliary coordinates.  In particular,
the theorem does not assert that every old quotient is divisible by \(f\),
or that the old and new equations have identical zero sets at unchanged
coordinates.  The local substitution proof also does not supply an
end-to-end Lean or Rocq certificate for the inherited native-window
compiler.

\subsection{Exact arithmetic size, degree, and companion equations}

The old evaluator already computes \(f^2\).  Its five private quotient
operations are replaced by the four operations
\[
 cT,\qquad cT-R,\qquad f^2(cT-R),\qquad f^2(cT-R)-c.
\]
Every other gate is retained.  The complete new evaluator has
\(38\) multiplications, \(15\) additions, and \(13\) subtractions.
Multiplication by a fixed numeral is charged; reuse is free.  All gates,
all fifteen witnesses, and all four coefficient ports are live.

The exact factor degrees in the order
\(N_0,N_1,N_a,N_k,N_t,N_{\mathrm{str}}\) are
\[
 22,\quad18,\quad62,\quad7,\quad2,\quad46.
\]
Put \(Q=\beta J\), \(k_0=\eta+\zeta\), and
\(C_1=Q-F-\alpha\).  The leading homogeneous polynomial of \(P_{66}\) is
\[
 -8Q^{91}h\gamma i^4k_0^{13}w^{13}s^{26}T^2f^4
       (wC_1-zQ).
\]
In particular, the coefficient of
\[
 J^{92}h\gamma i^4\eta^{13}w^{13}s^{26}T^2f^4z
\]
is \(8\beta^{92}\ne0\).  The exact degree is therefore 157, and the
final subtraction of \(\Delta\), of degree 12, cannot affect it.

For lower-degree alternatives, retain the new \(V\) but put
\[
 \widetilde N_a=\Delta(f^2-1)(V^2-y^2)+y^2,
 \qquad N_s=f^2-\Delta(ic^2)^2.
\]
The following three companion equations inherit the same positive
solvability equivalence:
\begin{align*}
 P_{68}&=N_0N_1\widetilde N_aN_kN_tN_s-1=0,\\
 P_{72}&=(N_0N_1-1)^2+(\widetilde N_aN_kN_t-1)^2+(N_s-1)^2=0,\\
 P_{74}&=(N_0-1)^2+(\widetilde N_a-1)^2+(N_s-1)^2
              +(N_1N_kN_t-1)^2=0.
\end{align*}
Indeed each is the corresponding inherited parent equation under the same
positive substitution \(T_{\mathrm{old}}=fT\); the established parent
equivalences and the reconstruction proof above apply.  Their fully paid
counts and exact degrees are
\[
\begin{array}{c|r|r|r|r|r}
 &\text{multiplications}&\text{additions}&\text{subtractions}
 &\text{total}&\text{degree}\\ \hline
 P_{66}&38&15&13&66&157\\
 P_{68}&39&15&14&68&113\\
 P_{72}&40&17&15&72&80\\
 P_{74}&40&18&16&74&68
\end{array}
\]
For \(P_{72}\), the first squared residual uniquely reaches degree 80;
its leading form is
\(64\gamma^2w^8s^{14}(\eta+\zeta)^6(\beta J)^{50}\).
For \(P_{74}\), the strong residual uniquely reaches degree 68 after
squaring; its leading form is
\(i^4w^4s^{12}(\eta+\zeta)^8(\beta J)^{40}\).
For \(P_{68}\), a degree-113 monomial has coefficient
\(8\beta^{64}\).  These nonzero coefficients certify the claimed exact
degrees uniformly for the prescribed tuple.

The supplied standard-library auditor checks the literal schedules, source
hashes, exact polynomial identities at the changed cut, leading homogeneous
forms, and full liveness.  Normal and optimized Python executions agree.
Its finite congruence examples and modular circuit diagnostics support the
implementation; the universal reconstruction lemma, rather than those
examples, supplies the new mathematical proof.  The result improves stated
repository tradeoffs and makes no claim of global minimality.


\section{Reproduction, certificates, and proof status}\label{sec:audit}
\subsection{What is fully explicit}
The archive supplies ordinary integer polynomials as finite, acyclic,
division-free arithmetic circuits. The data distinguishes input variables,
fixed integer leaves, and binary arithmetic nodes. All outputs depend on
the advertised variables, and every counted gate lies on the output's
dependency graph. There is no hidden existential parameter in a coefficient.

Both complete native tuples are included with their deterministic compilers
and literal machine/tile data. The low-variable equations use the coherent
package-31 tuple; the fifteen-variable equations use the coherent
package-36 tuple. Derived fixed leaves are reconstructed from explicit
integer recipes and checked against their payloads. Mathematical abbreviations
such as $\psi_A(R)$ occur only in proofs, never as polynomial gates.

The large fixed constants matter when interpreting simplicity.
These improvements concern the number of witnesses, total degree, and
the stated arithmetic-circuit cost. They do not claim a decrease in every
measure of description length or coefficient height. Expanding the complete
polynomials into lists of all monomials is unnecessary and would obscure
their exact finite definitions.

\subsection{Independent finite checks}
The supplied validators establish the following implementation facts.
\begin{enumerate}
\item \textbf{Exact source binding.}
An independently written sparse symbolic expansion agrees with the
parent arithmetic DAG over the polynomial ring in the native coefficient
ports. All 528 padded quadratic positions are distinct.
The 94 reversal records satisfy all 1,974 off-diagonal conditions.
Enumeration of all 49,632 convolution pairs gives exactly the intended
target coefficients, including terms absent from individual residuals.

\item \textbf{Literal finalizer replay.}
Independent evaluators reproduce the full ten-variable formula and the
nine-variable formula, coordinate translations, and cancellation-resolved
certificate. These are implementation diagnostics; their finite samples
are not represented as universal identities.

\item \textbf{Exact algebra and degrees.}
The three-root norm identity is compared symbolically as a 35-term
polynomial. The nine-variable telescoping identity is checked exactly.
Degree upper bounds and nonzero leading coefficients prove the reported
degrees. The short circuits have explicit exact cut identities and
leading homogeneous forms.

\item \textbf{Guard and toy checks.}
The separator audit includes 157,216 coprimality cases, 22,680 parity
cases, 7,750 Pell-remainder cases, and four explicit counterexamples to
dropping needed guards. Small accepting and rejecting carry-mask examples
exercise the source in nonvacuous situations. These tests supplement
the universal proofs above.

\item \textbf{Reconstruction and robust execution.}
The release checker rebuilds in a temporary copy and compares protected
artifacts byte for byte. Assertions needed for correctness are explicit
checks and remain active under Python's optimization flag.
\end{enumerate}

Run from the extracted archive:
\begin{verbatim}
python3 verify_all.py
\end{verbatim}
The mathematical verifiers use the Python standard library.
Building the article additionally requires a usual LaTeX installation:
\begin{verbatim}
cd report
pdflatex -interaction=nonstopmode -halt-on-error report.tex
pdflatex -interaction=nonstopmode -halt-on-error report.tex
\end{verbatim}
The README provides the exact file map and separate commands.

\subsection{What these certificates do and do not establish}
There are three distinct levels of justification.
First, the new selective-extraction, packed-successor, arithmetic-separator,
and squared-modulus results are proved as ordinary integer mathematics
in this report and its detailed companion proofs.
Second, the exact delivered circuits, coefficients, and degree/cost claims
have reproducible finite certificates.
Third, the semantic endpoint imports the repository's native-window
compiler theorem, fixed K215 calibration, and the existing Lucas/ratio
theorems identified in the provenance records.

The complete new chain has not been rebuilt in Lean or Rocq.
A finite circuit audit does not certify a machine's unbounded semantics,
and a randomized replay does not prove an algebraic identity.
The existing active Lean development proves several inherited local
identities and formal existence theorems; those are not silently promoted
to a kernel proof of the new explicit native-compiler composition.

All new reasoning uses finite algebra, elementary number theory, and
standard arithmetic theorems formalizable in ZFC. Subject to the stated
inherited mathematical theorems, the composition proves
\[
 \ZFC\vdash
 \bigl(\neg\exists\mathbf x\in D\;F(\mathbf x)=0\bigr)
 \longleftrightarrow\Con(\ZFC)
\]
for the specified domains $D$. This equivalence does not decide which
side holds. In particular, no experimental search for witnesses or
assumption of consistency was used to manufacture the coefficients.

\section{What changed, and promising next questions}
The improvements come from three structural changes and one separate
Pell refinement.

\paragraph{Use the support of each residual.}
Uniform block spacing protects many positions that no residual can use.
The new target rule forbids only actual cross-interference against
the full dense quadratic square. It reduces the maximum target from
$35\,177$ to $7\,243$ and the packed span from $36\,775$ to $8\,841$.
The target list is a finite certificate; the heuristic that found it
is not part of the proof.

\paragraph{Translate positivity after packing.}
The identity
$1+\sum_i(n_i+1)Z^{w_i}=1+\sum_i n_iZ^{w_i}+\ell(Z)$
avoids expanding a successor substitution separately in every residual.
Weighted records drop from 136 to 94 and port-expanded terms from 161
to 99. In these evaluators, the number of nonliteral fixed leaves drops
from 17 to 10. The small change of coefficient representation also
creates more freedom for target placement.

\paragraph{Separate radicals using arithmetic.}
The previous weights grow enough to dominate earlier square roots.
The new weights $(1,b,bK)$ exploit a coprime factorization,
the modulo-four structure of the last two radicands, and the constant
remainder in a Pell equation. This removes a large source of degree
without adding a quantified variable.
The unit case $b=1$ is essential: ignoring it would leave a gap
exactly where the positivity register vanishes.

\paragraph{Keep the full square modulus in the short equation.}
The stronger congruence modulo $f^2$ lets the reconstructed old quotient
be a multiple of $f$. Reparameterizing that quotient removes one paid
gate. This is a change in existential presentation, so its completeness
proof must reconstruct witnesses; a same-coordinate symbolic rewrite
alone would not suffice.

The highest-value next targets are consequently specific.
One is to jointly optimize the positive-parent presentation, Sidon
positions, target placement, and addition chains, with exact certificates
kept separate from search heuristics.
Another is to classify when a Pell remainder can replace a general
magnitude separator in other relation-combining constructions.
A third is to formalize the source interface and new separator as
reusable lemmas, then connect them to a fully checked native-compiler
endpoint. A lower variable count would require an additional structural
idea; none is claimed by the present work.

\section*{Attribution and integration notes}
\addcontentsline{toc}{section}{Attribution and integration notes}
The literal K215 machine, native-window compiler, original parents,
coefficient tuples, quintic merger, and earlier finalizers belong to the
cited ProveIt development and its credited mathematical sources.
Sun's ten-integer-variable theorem is not a new result of this report.
The nine-positive branch retains the Jones ratio theorem and the
package-45 strict-majorant finalizer.

The new contributions here are the support-sensitive target theorem and
certified schedule, the packed successor translation, the arithmetic
separator and its signed bootstrap, the squared-modulus quotient
reconstruction, and their explicit compiled combinations.
Exact gate sharing is identified separately from these mathematical changes.

The package is intended as a new research contribution beside the pinned
existing packages. It leaves those sources unchanged. A suggested directory
name is \file{46-sparse-arithmetic}; repository maintainers can choose the
final numbering. Preserve the included upstream attribution and licensing
notices when integrating copied coefficient and proof artifacts.

\phantomsection
\addcontentsline{toc}{section}{References}
\begin{thebibliography}{9}
\raggedright
\bibitem{repo}
V. Reshetnikov and contributors,
\emph{ProveIt}, source snapshot
\texttt{124c07c0b77a9e20e84676514d8b1ff78536ea3f},
reviewed 23 October 2026.
\url{https://github.com/VladimirReshetnikov/ProveIt}.
The complete file-level source pins are included in the archive.

\bibitem{repo40}
ProveIt, \emph{40-packing10}, in
\file{Computability/DiophantineConsistency/con-zfc-diophantine-forms},
at the source snapshot above.
The report, source parent, inherited finalizer proof, and notices
are bundled under \file{upstream/}.

\bibitem{repo45}
ProveIt, \emph{45-multioutput}, at the same source snapshot.
Its nine-variable proof and dependency attribution are bundled as
\file{upstream/PACKAGE45_NINE_PROOF.md}.

\bibitem{repo36}
ProveIt, \emph{36-small-tuple}, at the same source snapshot.
The full coefficient data and inherited P67 proof are bundled with
the short-equation contribution.

\bibitem{sun2021}
Z.-W. Sun, \emph{Further results on Hilbert's Tenth Problem},
Science China Mathematics \textbf{64} (2021), 281--306.
\url{https://arxiv.org/abs/1704.03504}.

\bibitem{sun2026}
Z.-W. Sun, \emph{Ten unknowns for Hilbert's tenth problem over the integers},
arXiv:2609.01594v2 (2026).
\url{https://arxiv.org/abs/2609.01594v2}.
The ten-integer composition uses the modified Lucas certificate,
residue interpolation, and binomial bridge, as identified in the
included inherited proof.
\end{thebibliography}
\end{document}
